\chapter{Generation Two: Robust Template Generation and Error Handling}
\label{chap:gen2}

\section{Overview}

Generation Two focuses on addressing critical operational challenges encountered in Generation One, particularly around template generation reliability, operator misuse, and data quality. Rather than introducing complex evolutionary algorithms, this generation prioritizes robust error handling, intelligent retry mechanisms, and systematic learning from failures to improve template generation success rates.

\begin{figure}[h]
\centering
\begin{tikzpicture}[
    node distance=1.5cm,
    gen1/.style={rectangle, draw=blue!50, fill=blue!10, thick, minimum width=3cm, minimum height=1.5cm, text centered, rounded corners},
    gen2/.style={rectangle, draw=green!50, fill=green!10, thick, minimum width=3cm, minimum height=1.5cm, text centered, rounded corners},
    arrow/.style={->, >=stealth, thick, blue!70},
    label/.style={text width=2.5cm, text centered, font=\small}
]
    % Generation One
    \node[gen1] (gen1) {Generation One\\\small AI Generation\\Concurrent Testing};
    
    % Generation Two additions
    \node[gen2, right=of gen1] (retry) {Retry\\Logic};
    \node[gen2, below=of gen1] (validation) {Template\\Validation};
    \node[gen2, below=of retry] (diversity) {Template\\Diversification};
    
    % Arrows showing evolution
    \draw[arrow] (gen1) -> node[above, label] {Enhanced with} (retry);
    \draw[arrow] (gen1) -> node[left, label] {Enhanced with} (validation);
    \draw[arrow] (retry) -> node[right, label] {Works with} (diversity);
    \draw[arrow] (validation) -> node[below, label] {Works with} (diversity);
    
    % Combined system
    \node[gen2, below=of diversity, xshift=-1.5cm, minimum width=6.5cm, fill=green!20] (gen2system) {Generation Two\\Robust System};
    
    \draw[arrow] (retry) -> (gen2system);
    \draw[arrow] (validation) -> (gen2system);
    \draw[arrow] (diversity) -> (gen2system);
\end{tikzpicture}
\caption{Evolution from Generation One to Generation Two}
\label{fig:gen1-to-gen2}
\end{figure}

\section{Key Improvements}

\subsection{Two-Phase Template Generation Architecture}

Generation Two implements a revolutionary two-phase approach that separates structural generation from operator/field selection:

\begin{figure}[h]
\centering
\begin{tikzpicture}[
    node distance=1.5cm and 2cm,
    phase1/.style={rectangle, draw=blue!50, fill=blue!10, thick, minimum width=4cm, minimum height=1.5cm, text centered, rounded corners},
    phase2/.style={rectangle, draw=green!50, fill=green!10, thick, minimum width=4cm, minimum height=1.5cm, text centered, rounded corners},
    step/.style={rectangle, draw=orange!50, fill=orange!10, thick, minimum width=3cm, minimum height=1cm, text centered},
    arrow/.style={->, >=stealth, thick},
    label/.style={text width=3cm, text centered, font=\small}
]
    % Phase 1: Algorithmic Generation
    \node[phase1] (phase1) {Phase 1: Algorithmic\\Placeholder Generation};
    \node[step, below=of phase1] (step4) {Step 4:\\Alpha Ideation};
    
    % Phase 2: Ollama Selection
    \node[phase2, right=of phase1] (phase2) {Phase 2: Ollama\\Index-Based Selection};
    \node[step, below=of phase2] (step5) {Step 5:\\Simulation};
    \node[step, below=of step5] (step6) {Step 6:\\Mining (20/80)};
    
    % Database
    \node[step, below=of step4, xshift=-1.5cm] (db) {Database\\Storage};
    
    % Arrows
    \draw[arrow] (phase1) -> (step4);
    \draw[arrow] (step4) -> node[left, label] {Store placeholders} (db);
    \draw[arrow] (phase2) -> (step5);
    \draw[arrow] (step5) -> (step6);
    \draw[arrow] (db) -> node[below, label] {80\% reuse} (step6);
    \draw[arrow] (step4) -> node[above, label] {Replace placeholders} (step5);
    \draw[arrow] (step6) -> node[right, label] {Replace placeholders} (step5);
\end{tikzpicture}
\caption{Two-Phase Template Generation Architecture}
\label{fig:two-phase-architecture}
\end{figure}

\subsection{Intelligent Retry Logic for Ollama Template Generation}

One of the primary challenges in Generation One was handling transient failures in Ollama API calls and template generation errors. Generation Two implements a sophisticated retry mechanism that learns from failures and adapts retry strategies.

\subsubsection{Exponential Backoff with Error Classification}

The retry system classifies errors and applies appropriate retry strategies:

\begin{lstlisting}[language=Python, caption=Intelligent Retry Handler]
class RetryHandler:
    def __init__(self, max_retries: int = 5, base_delay: float = 2.0):
        self.max_retries = max_retries
        self.base_delay = base_delay
        self.error_classification = {
            'transient': ['timeout', 'connection', 'rate_limit'],
            'permanent': ['syntax_error', 'invalid_operator', 'unknown_field'],
            'recoverable': ['missing_lookback', 'placeholder_field', 'event_input']
        }
        self.retry_stats = {
            'transient_retries': 0,
            'permanent_failures': 0,
            'recoverable_fixes': 0
        }
    
    def should_retry(self, error: Exception, attempt: int) -> Tuple[bool, float]:
        """Determine if error should be retried and calculate delay"""
        error_str = str(error).lower()
        
        # Classify error
        error_type = None
        for category, patterns in self.error_classification.items():
            if any(pattern in error_str for pattern in patterns):
                error_type = category
                break
        
        if error_type == 'permanent':
            # Don't retry permanent errors
            self.retry_stats['permanent_failures'] += 1
            return False, 0.0
        
        if error_type == 'recoverable':
            # Recoverable errors can be fixed, retry with longer delay
            self.retry_stats['recoverable_fixes'] += 1
            delay = self.base_delay * (2 ** attempt) * 1.5  # Extra delay for fixes
            return attempt < self.max_retries, delay
        
        if error_type == 'transient':
            # Transient errors: exponential backoff
            self.retry_stats['transient_retries'] += 1
            delay = self.base_delay * (2 ** attempt)
            return attempt < self.max_retries, delay
        
        # Unknown error: conservative retry
        delay = self.base_delay * (2 ** attempt)
        return attempt < self.max_retries, delay
\end{lstlisting}

\subsubsection{Template Generation with Retry}

The template generation process incorporates retry logic at multiple levels:

\begin{lstlisting}[language=Python, caption=Template Generation with Retry]
class OllamaTemplateGenerator:
    def generate_template_with_retry(
        self, 
        hypothesis: str,
        region: str,
        max_retries: int = 5
    ) -> Optional[str]:
        """Generate template with intelligent retry logic"""
        retry_handler = RetryHandler(max_retries=max_retries)
        
        for attempt in range(max_retries):
            try:
                # Generate template
                template = self.ollama_manager.generate_template(
                    hypothesis=hypothesis,
                    region=region,
                    available_operators=self.operators,
                    available_fields=self.get_fields_for_region(region)
                )
                
                if template:
                    # Validate template
                    is_valid, error_msg, suggested_fix = self.validator.validate_template(
                        template, region
                    )
                    
                    if is_valid:
                        return template
                    else:
                        # Validation failed - try to fix
                        if suggested_fix:
                            logger.info(f"🔧 Attempting fix: {suggested_fix}")
                            template = suggested_fix
                            # Re-validate fixed template
                            is_valid, _, _ = self.validator.validate_template(
                                template, region
                            )
                            if is_valid:
                                return template
                        
                        # Create error for retry logic
                        error = ValueError(f"Validation failed: {error_msg}")
                        should_retry, delay = retry_handler.should_retry(error, attempt)
                        
                        if not should_retry:
                            logger.error(f"❌ Template validation failed after {attempt + 1} attempts")
                            return None
                        
                        logger.info(f"🔄 Retrying template generation (attempt {attempt + 1}/{max_retries}) after {delay:.1f}s")
                        time.sleep(delay)
                        continue
                
                # No template generated
                error = ValueError("Ollama returned empty template")
                should_retry, delay = retry_handler.should_retry(error, attempt)
                if not should_retry:
                    return None
                time.sleep(delay)
                
            except requests.exceptions.Timeout as e:
                error = e
                should_retry, delay = retry_handler.should_retry(error, attempt)
                if not should_retry:
                    logger.error(f"❌ Timeout after {attempt + 1} attempts")
                    return None
                logger.warning(f"⏱️ Timeout, retrying in {delay:.1f}s...")
                time.sleep(delay)
                
            except requests.exceptions.ConnectionError as e:
                error = e
                should_retry, delay = retry_handler.should_retry(error, attempt)
                if not should_retry:
                    logger.error(f"❌ Connection error after {attempt + 1} attempts")
                    return None
                logger.warning(f"🔌 Connection error, retrying in {delay:.1f}s...")
                time.sleep(delay)
                
            except Exception as e:
                logger.error(f"❌ Unexpected error: {e}")
                error = e
                should_retry, delay = retry_handler.should_retry(error, attempt)
                if not should_retry:
                    return None
                time.sleep(delay)
        
        return None
\end{lstlisting}

\subsection{Common Ollama Template Generation Challenges}

Generation Two addresses several systematic issues where Ollama misuses operators from \texttt{operatorRAW.json}, leading to invalid templates.

\subsubsection{Missing Lookback Values}

Many time-series operators (e.g., \texttt{ts\_rank}, \texttt{ts\_mean}, \texttt{ts\_delta}) require a lookback parameter, but Ollama often generates templates without this parameter.

\begin{lstlisting}[language=Python, caption=Missing Lookback Detection and Fix]
class TemplateValidator:
    def _fix_missing_lookback(self, template: str, error_message: str) -> Tuple[str, List[str]]:
        """Fix missing lookback parameter by adding default lookback value"""
        import re
        fixes = []
        fixed_template = template
        
        # Operators that require lookback
        lookback_operators = [
            'ts_rank', 'ts_sum', 'ts_mean', 'ts_max', 'ts_min', 'ts_std',
            'ts_delta', 'ts_correlation', 'ts_covariance', 'ts_decay',
            'ts_product', 'ts_argmax', 'ts_argmin', 'ts_scale', 'ts_winsorize',
            'ts_zscore', 'ts_log', 'ts_reverse', 'ts_abs'
        ]
        
        for op in lookback_operators:
            # Pattern: operator(field) but not operator(field, ...)
            pattern = r'\b' + re.escape(op) + r'\s*\(\s*([^,)]+)\s*\)(?!\s*,\s*\d)'
            matches = list(re.finditer(pattern, fixed_template, re.IGNORECASE))
            
            if matches:
                # Replace from end to start to preserve positions
                for match in reversed(matches):
                    field_expr = match.group(1)
                    # Add lookback parameter: operator(field) -> operator(field, 20)
                    replacement = f"{op}({field_expr}, 20)"
                    fixed_template = (
                        fixed_template[:match.start()] + 
                        replacement + 
                        fixed_template[match.end():]
                    )
                    fixes.append(f"Added lookback parameter (20) to {op}")
                    logger.info(f"🔧 Added lookback to {op}: {op}({field_expr}) -> {op}({field_expr}, 20)")
        
        return fixed_template, fixes
\end{lstlisting}

\subsection{Algorithmic Placeholder Generation and Index-Based Selection}

Generation Two implements a two-phase template generation approach that separates structural generation from operator/field selection, significantly improving reliability and reducing errors. The system includes comprehensive validation to ensure all placeholders are replaced before submission, preventing errors from unreplaced placeholders.

\subsubsection{Phase 1: Algorithmic Placeholder Generation (Step 4)}

In Step 4 (Alpha Ideation and Generation), templates are generated algorithmically without calling Ollama. The system uses mathematical algorithms (random walk, Brownian motion, tree structures, linear chains) to create placeholder expressions:

\begin{lstlisting}[language=Python, caption=Algorithmic Placeholder Generation]
class AlgorithmicTemplateGenerator:
    """Generates placeholder expressions using mathematical algorithms"""
    
    def generate_placeholder_expression(
        self, 
        max_operators: int = 5,
        method: str = 'random'
    ) -> str:
        """
        Generate placeholder expression structure
        
        Methods:
        - 'random': Random walk through expression space
        - 'brownian': Brownian motion-inspired generation
        - 'tree': Binary tree structure
        - 'linear': Linear chain of operators
        """
        if method == 'random':
            return self._random_walk_generation(max_operators)
        elif method == 'brownian':
            return self._brownian_motion_generation(max_operators)
        elif method == 'tree':
            return self._tree_generation(max_operators)
        elif method == 'linear':
            return self._linear_generation(max_operators)
        
        # Example output: "OPERATOR1(OPERATOR2(DATA_FIELD1), 5)"
\end{lstlisting}

Key advantages of algorithmic generation:
\begin{itemize}
    \item \textbf{No LLM dependency}: Generation is deterministic and fast
    \item \textbf{Structural correctness}: Algorithms enforce max operator limits and avoid consecutive duplicates
    \item \textbf{Placeholder format}: Templates use \texttt{OPERATOR1}, \texttt{OPERATOR2}, \texttt{DATA\_FIELD1}, etc.
    \item \textbf{Database storage}: Placeholder templates are stored for later reuse
    \item \textbf{Duplicate checking}: System checks against database before storing new templates
\end{itemize}

\subsubsection{Phase 2: Ollama Index-Based Selection with Validation (Step 5 \& Step 6)}

In Step 5 (Simulation and Testing) and Step 6 (Continuous Mining), placeholders are replaced using Ollama's index-based selection with comprehensive validation to ensure \textbf{ALL} placeholders are replaced before submission.

The system presents Ollama with:

\begin{enumerate}
    \item The placeholder template (e.g., \texttt{OPERATOR1(OPERATOR2(DATA\_FIELD1), 5)})
    \item \textbf{Explicit list of ALL required placeholders} found in the template
    \item A randomly shuffled list of operators with full metadata:
    \begin{itemize}
        \item Operator name and category
        \item Number of inputs required (including minimum inputs)
        \item Scope (REGULAR, MATRIX, VECTOR, COMBO, SELECTION)
        \item Event input support (Yes/No)
        \item Lookback usage (Yes/No)
        \item Description and definition
        \item Example usage
    \end{itemize}
    \item A randomly shuffled list of data fields (up to 50 for diversity)
\end{enumerate}

Ollama returns JSON with index selections, and the system validates completeness:

\begin{lstlisting}[language=Python, caption=Ollama Index-Based Selection with Validation]
def replace_placeholders_with_selection(
    placeholder_expression: str,
    available_operators: List[Dict],
    available_fields: List[Dict]
) -> Optional[str]:
    """
    Replace placeholders with validation
    
    Returns:
    - Replaced template if ALL placeholders replaced
    - None if any placeholders remain (indicates failure)
    """
    # 1. Discover all placeholders in template (case-insensitive)
    operator_placeholders = find_all_operator_placeholders(template)
    field_placeholders = find_all_field_placeholders(template)
    
    # 2. Build prompt with explicit list of required placeholders
    prompt = build_selection_prompt(
        placeholder_expression,
        required_operators=operator_placeholders,  # Explicit list
        required_fields=field_placeholders,        # Explicit list
        shuffled_operators=shuffled_operators,
        shuffled_fields=shuffled_fields
    )
    
    # 3. Get Ollama selection
    selection = ollama.generate(prompt)
    
    # 4. Filter out invalid selections (placeholders not in template)
    selection = filter_valid_selections(selection, operator_placeholders, field_placeholders)
    
    # 5. Replace placeholders mechanically
    result = replace_placeholders_mechanically(
        placeholder_expression,
        selection,
        operator_mapping,
        field_mapping
    )
    
    # 6. CRITICAL: Verify ALL placeholders were replaced
    remaining = find_remaining_placeholders(result)
    if remaining:
        logger.error(f"Placeholders not replaced: {remaining}")
        return None  # Indicate failure
    
    return result
\end{lstlisting}

\paragraph{Validation and Retry Logic}

Before submission, the system performs multiple validation checks:

\begin{lstlisting}[language=Python, caption=Pre-Submission Validation]
def run_simulation(template, region):
    # 1. Replace placeholders
    replaced = replace_placeholders_with_selection(template, ...)
    
    if replaced:
        template = replaced
        # 2. Verify all placeholders replaced
        remaining = find_remaining_placeholders(template)
        
        if remaining:
            # 3. Retry replacement once
            replaced_retry = replace_placeholders_with_selection(template, ...)
            if replaced_retry:
                template = replaced_retry
                remaining = find_remaining_placeholders(template)
                
                if remaining:
                    # 4. FAIL: Skip submission
                    mark_slot_failed(f"Placeholders remain: {remaining}")
                    return
            else:
                # Retry failed
                mark_slot_failed("Placeholder replacement failed")
                return
    
    # 5. FINAL CHECK before submission
    remaining = find_remaining_placeholders(template)
    if remaining:
        mark_slot_failed(f"Cannot submit: placeholders remain")
        return
    
    # 6. Submit only if ALL placeholders replaced
    submit_simulation(template)
\end{lstlisting}

Benefits of validated index-based selection:
\begin{itemize}
    \item \textbf{No spelling errors}: Ollama selects by index, not by generating names
    \item \textbf{Rich metadata}: Full operator/field information guides selection
    \item \textbf{Random shuffling}: Ensures diversity across selections
    \item \textbf{Mechanical replacement}: Guaranteed correct replacement using mappings
    \item \textbf{Complete validation}: System verifies ALL placeholders are replaced
    \item \textbf{Automatic retry}: Retries replacement if incomplete
    \item \textbf{Submission safety}: Never submits templates with unreplaced placeholders
    \item \textbf{Fallback support}: Falls back to old method if Ollama selection fails
\end{itemize}

\subsubsection{Step 6: Continuous Mining with 20/80 Logic}

Step 6 (Continuous Alpha Mining) implements an intelligent template reuse strategy:

\begin{itemize}
    \item \textbf{20\% chance}: Generate new placeholder template algorithmically (different from database)
    \item \textbf{80\% chance}: Reuse unsimulated template from database
    \item \textbf{Placeholder replacement with validation}: All templates (new or reused) go through Ollama index-based selection with complete validation before simulation
    \item \textbf{Duplicate prevention}: System checks templates aren't currently being simulated
    \item \textbf{Database tracking}: Tracks which templates have been simulated
    \item \textbf{Placeholder normalization}: Deduplication treats templates with permuted placeholder numbers as equivalent (e.g., \texttt{OPERATOR3(OPERATOR2(...))} = \texttt{OPERATOR1(OPERATOR2(...))})
\end{itemize}

\begin{lstlisting}[language=Python, caption=Step 6 Mining Logic with Validation]
def _process_simulations(self):
    """Process simulations with 20/80 logic and validation"""
    for each_slot:
        if random() < 0.2:  # 20% chance
            # Generate new placeholder template
            template = generate_new_placeholder_template()
            # Ensure different from database
            while template in database:
                template = generate_new_placeholder_template()
            store_in_database(template)
        else:  # 80% chance
            # Pick unsimulated template from database
            template = pick_unsimulated_template_from_db()
            if not template:
                # Fallback: generate new
                template = generate_new_placeholder_template()
        
        # Replace placeholders using Ollama selection with validation
        replaced = replace_placeholders_with_selection(template, ...)
        if not replaced:
            # Replacement failed, skip this template
            mark_slot_failed("Placeholder replacement failed")
            continue
        
        template = replaced
        
        # Verify all placeholders replaced
        remaining = find_remaining_placeholders(template)
        if remaining:
            # Retry once
            replaced_retry = replace_placeholders_with_selection(template, ...)
            if replaced_retry:
                template = replaced_retry
                remaining = find_remaining_placeholders(template)
                if remaining:
                    mark_slot_failed(f"Placeholders remain: {remaining}")
                    continue
            else:
                mark_slot_failed("Placeholder replacement failed")
                continue
        
        # Final check before submission
        remaining = find_remaining_placeholders(template)
        if remaining:
            mark_slot_failed(f"Cannot submit: placeholders remain")
            continue
        
        # Submit only if ALL placeholders replaced
        submit_simulation(template)
\end{lstlisting}

\subsubsection{Refeed Mechanism with Placeholder Replacement and Validation}

When templates fail during simulation, the refeed mechanism uses Ollama index-based selection with the same validation logic:

\begin{lstlisting}[language=Python, caption=Refeed with Placeholder Replacement and Validation]
def handle_refeed(template, error_message):
    # 1. Check if template still has placeholders BEFORE fixing
    if has_placeholders(template):
        # Replace placeholders BEFORE fixing
        template = replace_placeholders_with_selection(template)
        if not template:  # Replacement failed
            return None

    # 2. Apply refeed correction
    fixed_template = refeed_with_correction(template, error_message)

    # 3. Check if fixed template reintroduced placeholders
    if has_placeholders(fixed_template):
        # Replace again
        fixed_template = replace_placeholders_with_selection(fixed_template)
        if not fixed_template:  # Replacement failed
            return None

    # 4. FINAL VALIDATION before resubmission
    remaining = find_remaining_placeholders(fixed_template)
    if remaining:
        logger.error(f"Cannot resubmit: placeholders remain: {remaining}")
        return None  # Cannot proceed

    # 5. Retry simulation only if ALL placeholders replaced
    return retry_simulation(fixed_template)
\end{lstlisting}

This ensures that:
\begin{itemize}
    \item Placeholders are always replaced before any fix attempts
    \item No \texttt{"operator1"} or \texttt{"operator4"} errors occur from unreplaced placeholders
    \item All templates submitted to API are fully replaced (no placeholders)
    \item Refeed mechanism validates completeness before resubmission
    \item System fails gracefully if replacement cannot complete
    \item Both Step 5 and Step 6 use identical validation logic
\end{itemize}

\subsubsection{Placeholder Normalization for Deduplication}

To ensure deduplication is not affected by placeholder number permutations, the system normalizes placeholders before comparison:

\begin{lstlisting}[language=Python, caption=Placeholder Normalization]
def normalize_placeholders(template: str) -> str:
    """
    Normalize placeholder numbers to canonical form
    
    Example:
    - OPERATOR3(OPERATOR2(OPERATOR1(DATA_FIELD1))) 
    - OPERATOR1(OPERATOR2(OPERATOR3(DATA_FIELD1)))
    Both become: OPERATOR_A(OPERATOR_B(OPERATOR_C(DATA_FIELD_A))
    """
    # Find all placeholders in order of appearance
    operator_placeholders = find_all_operator_placeholders(template)
    field_placeholders = find_all_field_placeholders(template)
    
    # Replace by occurrence order (left to right)
    # First occurrence -> OPERATOR_A, second -> OPERATOR_B, etc.
    for idx, placeholder in enumerate(operator_placeholders):
        canonical = f"OPERATOR_{chr(65 + idx)}"  # A, B, C, ...
        template = replace(placeholder, canonical)
    
    # Same for fields
    for idx, placeholder in enumerate(field_placeholders):
        canonical = f"DATA_FIELD_{chr(65 + idx)}"  # A, B, C, ...
        template = replace(placeholder, canonical)
    
    return template
\end{lstlisting}

This ensures that templates with different placeholder numberings are treated as equivalent during duplicate checking, preventing false duplicates from being generated.

\subsubsection{Legacy: Placeholder Data Fields Instead of Actual Fields}

The previous approach (now used as fallback) replaced placeholders mechanically without Ollama selection. This is still available as a fallback mechanism when Ollama selection fails.

\begin{lstlisting}[language=Python, caption=Placeholder Field Replacement]
class TemplateValidator:
    def _replace_placeholder_fields(
        self, 
        template: str, 
        region: str, 
        delay: int = None
    ) -> Tuple[str, List[str]]:
        """Replace placeholder fields with actual data fields"""
        import re
        fixes = []
        fixed_template = template
        
        # Detect placeholder patterns
        placeholder_pattern = r'\b(DATA_FIELD\d+|data_field\d+|FIELD\d+|field\d+)\b'
        placeholders = re.findall(placeholder_pattern, template, re.IGNORECASE)
        
        if not placeholders:
            return template, fixes
        
        # Get valid fields for region
        valid_fields = self.get_data_fields_for_region(region, delay)
        if not valid_fields:
            logger.warning(f"⚠️ No valid fields available for region {region}")
            return template, fixes
        
        # Map placeholders to actual fields (round-robin to ensure diversity)
        placeholder_to_field = {}
        field_index = 0
        
        for placeholder in set(placeholders):
            # Use different fields for different placeholders
            field = valid_fields[field_index % len(valid_fields)]
            placeholder_to_field[placeholder] = field['id']
            field_index += 1
        
        # Replace placeholders
        for placeholder, field_id in placeholder_to_field.items():
            pattern = r'\b' + re.escape(placeholder) + r'\b'
            fixed_template = re.sub(
                pattern, 
                field_id, 
                fixed_template, 
                flags=re.IGNORECASE
            )
            fixes.append(f"Replaced {placeholder} with {field_id}")
            logger.info(f"🔧 Replaced placeholder: {placeholder} -> {field_id}")
        
        return fixed_template, fixes
\end{lstlisting}

\subsubsection{Event Input Data Field Incompatibility}

Certain operators do not accept event input data fields, but Ollama may generate templates that use event inputs with incompatible operators. Generation Two tracks operator-field compatibility pairs and learns from simulation errors.

\begin{lstlisting}[language=Python, caption=Event Input Compatibility Tracking]
class EventInputCompatibilityTracker:
    def __init__(self, db_path: str = "generation_two_backtests.db"):
        self.db_path = db_path
        self._init_database()
    
    def _init_database(self):
        """Initialize database tables for compatibility tracking"""
        import sqlite3
        conn = sqlite3.connect(self.db_path)
        cursor = conn.cursor()
        
        # Table for event input fields
        cursor.execute('''
            CREATE TABLE IF NOT EXISTS event_input_fields (
                field_id TEXT PRIMARY KEY,
                region TEXT,
                delay INTEGER,
                is_event_input INTEGER,
                first_detected TIMESTAMP DEFAULT CURRENT_TIMESTAMP,
                last_updated TIMESTAMP DEFAULT CURRENT_TIMESTAMP
            )
        ''')
        
        # Table for incompatible operator-field pairs
        cursor.execute('''
            CREATE TABLE IF NOT EXISTS incompatible_pairs (
                operator_name TEXT,
                field_id TEXT,
                region TEXT,
                error_message TEXT,
                first_encountered TIMESTAMP DEFAULT CURRENT_TIMESTAMP,
                occurrence_count INTEGER DEFAULT 1,
                PRIMARY KEY (operator_name, field_id, region)
            )
        ''')
        
        conn.commit()
        conn.close()
    
    def record_incompatibility(
        self, 
        operator_name: str, 
        field_id: str, 
        region: str,
        error_message: str
    ):
        """Record an incompatible operator-field pair"""
        import sqlite3
        conn = sqlite3.connect(self.db_path)
        cursor = conn.cursor()
        
        # Check if pair already exists
        cursor.execute('''
            SELECT occurrence_count FROM incompatible_pairs
            WHERE operator_name = ? AND field_id = ? AND region = ?
        ''', (operator_name, field_id, region))
        
        result = cursor.fetchone()
        if result:
            # Update occurrence count
            cursor.execute('''
                UPDATE incompatible_pairs
                SET occurrence_count = occurrence_count + 1,
                    error_message = ?
                WHERE operator_name = ? AND field_id = ? AND region = ?
            ''', (error_message, operator_name, field_id, region))
        else:
            # Insert new pair
            cursor.execute('''
                INSERT INTO incompatible_pairs 
                (operator_name, field_id, region, error_message)
                VALUES (?, ?, ?, ?)
            ''', (operator_name, field_id, region, error_message))
        
        conn.commit()
        conn.close()
    
    def is_incompatible(
        self, 
        operator_name: str, 
        field_id: str, 
        region: str
    ) -> bool:
        """Check if operator-field pair is known to be incompatible"""
        import sqlite3
        conn = sqlite3.connect(self.db_path)
        cursor = conn.cursor()
        
        cursor.execute('''
            SELECT 1 FROM incompatible_pairs
            WHERE operator_name = ? AND field_id = ? AND region = ?
        ''', (operator_name, field_id, region))
        
        result = cursor.fetchone()
        conn.close()
        return result is not None
    
    def mark_as_event_input(self, field_id: str, region: str, delay: int = None):
        """Mark a field as an event input"""
        import sqlite3
        conn = sqlite3.connect(self.db_path)
        cursor = conn.cursor()
        
        cursor.execute('''
            INSERT OR REPLACE INTO event_input_fields 
            (field_id, region, delay, is_event_input, last_updated)
            VALUES (?, ?, ?, 1, CURRENT_TIMESTAMP)
        ''', (field_id, region, delay))
        
        conn.commit()
        conn.close()
\end{lstlisting}

\subsection{Template Diversification}

Generation Two implements template diversification at two levels: (1) across concurrent generation threads, and (2) within the simulated alpha pool to avoid redundant testing.

\subsubsection{Cross-Thread Diversification}

Each generation thread maintains awareness of templates generated by other threads to avoid duplicates:

\begin{lstlisting}[language=Python, caption=Cross-Thread Template Diversification]
class TemplateDiversifier:
    def __init__(self):
        self.generated_templates = set()  # Thread-safe set
        self.operator_usage = {}  # Track operator usage across threads
        self.lock = threading.Lock()
    
    def is_duplicate(self, template: str) -> bool:
        """Check if template was already generated"""
        with self.lock:
            # Normalize template for comparison
            normalized = self._normalize_template(template)
            return normalized in self.generated_templates
    
    def register_template(self, template: str):
        """Register a generated template"""
        with self.lock:
            normalized = self._normalize_template(template)
            self.generated_templates.add(normalized)
            
            # Track operator usage
            operators = self._extract_operators(template)
            for op in operators:
                self.operator_usage[op] = self.operator_usage.get(op, 0) + 1
    
    def get_diverse_operators(
        self, 
        available_operators: List[Dict],
        num_operators: int = 3
    ) -> List[Dict]:
        """Select diverse operators (prioritize less-used operators)"""
        with self.lock:
            # Sort by usage count (ascending)
            sorted_ops = sorted(
                available_operators,
                key=lambda op: self.operator_usage.get(op['name'], 0)
            )
            
            # Select diverse operators
            selected = []
            for op in sorted_ops:
                if len(selected) >= num_operators:
                    break
                if op['name'] not in [s['name'] for s in selected]:
                    selected.append(op)
            
            return selected
\end{lstlisting}

\subsubsection{Simulated Alpha Pool Diversification}

Before submitting templates for simulation, Generation Two checks against the pool of already-simulated alphas to avoid redundant testing:

\begin{lstlisting}[language=Python, caption=Alpha Pool Diversification]
class AlphaPoolDiversifier:
    def __init__(self, db_path: str = "generation_two_backtests.db"):
        self.db_path = db_path
        self._init_database()
    
    def _init_database(self):
        """Initialize database for alpha pool tracking"""
        import sqlite3
        conn = sqlite3.connect(self.db_path)
        cursor = conn.cursor()
        
        cursor.execute('''
            CREATE TABLE IF NOT EXISTS simulated_alphas (
                template_hash TEXT PRIMARY KEY,
                template TEXT,
                region TEXT,
                delay INTEGER,
                sharpe REAL,
                fitness REAL,
                simulated_at TIMESTAMP DEFAULT CURRENT_TIMESTAMP
            )
        ''')
        
        # Index for fast lookups
        cursor.execute('''
            CREATE INDEX IF NOT EXISTS idx_template_hash 
            ON simulated_alphas(template_hash)
        ''')
        
        conn.commit()
        conn.close()
    
    def has_been_simulated(self, template: str, region: str) -> bool:
        """Check if template has already been simulated"""
        import sqlite3
        import hashlib
        
        template_hash = self._hash_template(template, region)
        
        conn = sqlite3.connect(self.db_path)
        cursor = conn.cursor()
        
        cursor.execute('''
            SELECT 1 FROM simulated_alphas
            WHERE template_hash = ?
        ''', (template_hash,))
        
        result = cursor.fetchone()
        conn.close()
        
        return result is not None
    
    def register_simulation(
        self, 
        template: str, 
        region: str, 
        delay: int,
        sharpe: float = None,
        fitness: float = None
    ):
        """Register a simulated alpha"""
        import sqlite3
        template_hash = self._hash_template(template, region)
        
        conn = sqlite3.connect(self.db_path)
        cursor = conn.cursor()
        
        cursor.execute('''
            INSERT OR REPLACE INTO simulated_alphas
            (template_hash, template, region, delay, sharpe, fitness)
            VALUES (?, ?, ?, ?, ?, ?)
        ''', (template_hash, template, region, delay, sharpe, fitness))
        
        conn.commit()
        conn.close()
    
    def _hash_template(self, template: str, region: str) -> str:
        """Generate hash for template"""
        import hashlib
        normalized = self._normalize_template(template)
        combined = f"{normalized}:{region}"
        return hashlib.sha256(combined.encode()).hexdigest()
    
    def _normalize_template(self, template: str) -> str:
        """Normalize template for comparison"""
        import re
        # Remove whitespace
        normalized = re.sub(r'\s+', '', template)
        # Convert to lowercase
        normalized = normalized.lower()
        return normalized
\end{lstlisting}

\subsection{Production Correlation Loading}

Generation One encountered API limitations when loading production correlation data. Generation Two implements a robust correlation loading system with rate limiting, polling, and error handling.

\begin{lstlisting}[language=Python, caption=Production Correlation Loader with API Handling]
class ProductionCorrelationLoader:
    def __init__(self, credentials_path: str, max_retries: int = 10):
        self.credentials_path = credentials_path
        self.max_retries = max_retries
        self.rate_limiter = RateLimiter(requests_per_minute=30)
        self.setup_auth()
    
    def load_correlation_data(
        self, 
        alpha_id: str,
        poll_interval: int = 5,
        max_polls: int = 10
    ) -> Optional[Dict]:
        """Load correlation data with polling for async processing"""
        url = f"https://api.worldquantbrain.com/alphas/{alpha_id}/correlations/prod"
        
        for poll_attempt in range(max_polls):
            try:
                # Rate limiting
                self.rate_limiter.wait_if_needed()
                
                response = self.make_api_request('GET', url)
                
                # Handle rate limiting (429)
                if response.status_code == 429:
                    retry_after = float(response.headers.get('Retry-After', 2.0))
                    wait_time = retry_after * (2 ** poll_attempt)
                    logger.warning(f"⏱️ Rate limited, waiting {wait_time:.1f}s...")
                    time.sleep(wait_time)
                    continue
                
                if response.status_code == 200:
                    corr_data = response.json()
                    
                    # Check if we have actual data (not empty response while processing)
                    has_data = (
                        corr_data.get('max') is not None or 
                        corr_data.get('min') is not None or 
                        corr_data.get('records') or
                        len(corr_data) > 0
                    )
                    
                    if has_data:
                        logger.info(f"✅ Correlation data loaded for alpha {alpha_id}")
                        return corr_data
                    else:
                        # Data still processing, poll again
                        if poll_attempt < max_polls - 1:
                            logger.info(
                                f"⏳ Correlation data still processing "
                                f"(attempt {poll_attempt + 1}/{max_polls}), "
                                f"waiting {poll_interval}s..."
                            )
                            time.sleep(poll_interval)
                            continue
                        else:
                            logger.warning(
                                f"⚠️ Correlation data not ready after "
                                f"{max_polls} polling attempts"
                            )
                            return None
                
                elif response.status_code == 404:
                    logger.warning(f"⚠️ Alpha {alpha_id} not found")
                    return None
                else:
                    logger.error(
                        f"❌ Unexpected status code {response.status_code}: "
                        f"{response.text[:200]}"
                    )
                    return None
                    
            except requests.exceptions.Timeout:
                logger.warning(f"⏱️ Timeout loading correlation data (attempt {poll_attempt + 1})")
                if poll_attempt < max_polls - 1:
                    time.sleep(poll_interval)
                    continue
                return None
                
            except Exception as e:
                logger.error(f"❌ Error loading correlation data: {e}")
                if poll_attempt < max_polls - 1:
                    time.sleep(poll_interval)
                    continue
                return None
        
        return None
\end{lstlisting}

\subsection{Flatline PnL Handling}

Generation One (v2) introduced flatline PnL detection, where alphas produce constant or near-constant PnL values over time. While this is tracked for monitoring purposes, Generation Two does not incorporate flatline PnL into the reward function, as the current focus is on breadth-first search and learning from all mined simulations.

\begin{lstlisting}[language=Python, caption=Flatline PnL Detection (Monitoring Only)]
class FlatlinePnLMonitor:
    def __init__(self):
        self.flatline_stats = {
            'detected': 0,
            'total_checked': 0
        }
    
    def detect_flatline(self, pnl_values: List[float]) -> bool:
        """Detect flatline PnL pattern (for monitoring only)"""
        if len(pnl_values) < 10:
            return False
        
        self.flatline_stats['total_checked'] += 1
        
        # Calculate statistics
        pnl_min = min(pnl_values)
        pnl_max = max(pnl_values)
        pnl_range = pnl_max - pnl_min
        pnl_mean = sum(pnl_values) / len(pnl_values)
        pnl_std = (
            sum((x - pnl_mean) ** 2 for x in pnl_values) / len(pnl_values)
        ) ** 0.5
        
        # Flatline detection criteria
        is_flatline = (
            pnl_range < 0.001 or  # Very small range
            pnl_std < 0.0001 or   # Very low standard deviation
            (pnl_max - pnl_min) < 0.0001  # Very small variance
        )
        
        if is_flatline:
            self.flatline_stats['detected'] += 1
            logger.warning(
                f"🔴 Flatline PnL detected: range={pnl_range:.6f}, "
                f"std={pnl_std:.6f}"
            )
        
        return is_flatline
    
    def get_stats(self) -> Dict:
        """Get flatline detection statistics"""
        return {
            'detected': self.flatline_stats['detected'],
            'total_checked': self.flatline_stats['total_checked'],
            'detection_rate': (
                self.flatline_stats['detected'] / 
                max(1, self.flatline_stats['total_checked'])
            )
        }
\end{lstlisting}

\textbf{Note}: Flatline PnL detection is used for monitoring and logging purposes only. It does not affect the reward function or template selection, as Generation Two prioritizes breadth-first exploration to learn from all simulation outcomes.

\section{Expected Improvements}

\subsection{Performance Metrics}

Generation Two is expected to achieve:

\begin{table}[h]
\centering
\begin{tabular}{lcc}
\toprule
Metric & Generation One & Generation Two (Expected) \\
\midrule
Template Generation Success Rate & 60-70\% & 85-95\% \\
Retry Success Rate & N/A & 40-50\% of initial failures \\
Validation Fix Rate & N/A & 60-70\% of invalid templates \\
Duplicate Template Rate & 15-20\% & <5\% \\
Correlation Loading Success Rate & 70-80\% & 90-95\% \\
\bottomrule
\end{tabular}
\caption{Expected Performance Improvements}
\end{table}

\subsection{Key Advantages}

\begin{enumerate}
    \item \textbf{Robust Error Handling}: Intelligent retry logic recovers from transient failures
    \item \textbf{Automatic Template Fixing}: Common operator misuse patterns are automatically corrected
    \item \textbf{Event Input Learning}: System learns which operators are incompatible with event inputs
    \item \textbf{Template Diversification}: Reduces redundant testing across threads and alpha pool
    \item \textbf{API Resilience}: Handles rate limiting and async processing for correlation data
\end{enumerate}

\section{Implementation Strategy}

\subsection{Integration with Generation One}

Generation Two extends Generation One's template generation pipeline with validation and retry layers:

\begin{lstlisting}[language=Python, caption=Generation Two Integration]
class EnhancedTemplateGeneratorV3(EnhancedTemplateGeneratorV2):
    """Generation Two: Extends Generation One with robust error handling"""
    
    def __init__(self, *args, **kwargs):
        super().__init__(*args, **kwargs)
        
        # Add Generation Two components
        self.retry_handler = RetryHandler(max_retries=5)
        self.validator = TemplateValidator(
            operators=self.operators,
            data_fields=self.data_fields,
            ollama_manager=self.ollama_manager,
            use_ast=True  # Enable AST-based validation
        )
        self.diversifier = TemplateDiversifier()
        self.alpha_pool_diversifier = AlphaPoolDiversifier()
        self.event_tracker = EventInputCompatibilityTracker()
        self.correlation_loader = ProductionCorrelationLoader(
            credentials_path=self.credentials_path
        )
        self.flatline_monitor = FlatlinePnLMonitor()
        
    def generate_template_robust(
        self, 
        hypothesis: str,
        region: str,
        delay: int = None
    ) -> Optional[str]:
        """Generate template with robust error handling"""
        # Check for duplicates
        if self.diversifier.is_duplicate(hypothesis):
            logger.debug(f"⚠️ Duplicate hypothesis detected, skipping")
            return None
        
        # Generate with retry
        template = self.generate_template_with_retry(
            hypothesis=hypothesis,
            region=region,
            max_retries=5
        )
        
        if not template:
            return None
        
        # Validate and fix
        is_valid, error_msg, suggested_fix = self.validator.validate_template(
            template, region, delay
        )
        
        if not is_valid:
            # Try to fix
            fixed_template, fixes = self.validator.fix_template(
                template, error_msg, region, delay
            )
            
            # Re-validate
            is_valid, _, _ = self.validator.validate_template(
                fixed_template, region, delay
            )
            
            if is_valid:
                template = fixed_template
                logger.info(f"✅ Template fixed: {fixes}")
            else:
                logger.warning(f"⚠️ Could not fix template: {error_msg}")
                return None
        
        # Register template
        self.diversifier.register_template(template)
        
        return template
    
    def simulate_template_robust(
        self,
        template: str,
        region: str,
        delay: int = None
    ) -> Optional[AlphaResult]:
        """Simulate template with duplicate checking"""
        # Check if already simulated
        if self.alpha_pool_diversifier.has_been_simulated(template, region):
            logger.debug(f"⚠️ Template already simulated, skipping")
            return None
        
        # Simulate
        result = super().simulate_template(template, region, delay)
        
        if result:
            # Register simulation
            self.alpha_pool_diversifier.register_simulation(
                template=template,
                region=region,
                delay=delay,
                sharpe=result.sharpe if hasattr(result, 'sharpe') else None,
                fitness=result.fitness if hasattr(result, 'fitness') else None
            )
            
            # Monitor flatline PnL (for logging only)
            if hasattr(result, 'pnl_values'):
                is_flatline = self.flatline_monitor.detect_flatline(result.pnl_values)
                if is_flatline:
                    logger.warning(f"🔴 Flatline PnL detected for template")
        
        return result
\end{lstlisting}

\section{Challenges and Solutions}

\subsection{Operator Misuse Patterns}

\textbf{Challenge}: Ollama frequently misuses operators from \texttt{operatorRAW.json}, generating invalid templates.

\textbf{Solution}:
\begin{itemize}
    \item Pattern-based detection and fixing (missing lookback, placeholder fields)
    \item AST-based validation before simulation
    \item Learning from simulation errors to build compatibility database
    \item Enhanced prompts that emphasize operator requirements
\end{itemize}

\subsection{Event Input Incompatibility}

\textbf{Challenge}: Not all operators accept event input data fields, but this information is not readily available.

\textbf{Solution}:
\begin{itemize}
    \item Track operator-field pairs that cause simulation errors
    \item Build database of incompatible pairs over time
    \item Use learned knowledge to prevent invalid combinations
    \item Mark fields as event inputs when detected
\end{itemize}

\subsection{API Rate Limiting}

\textbf{Challenge}: Production correlation API has rate limits and async processing.

\textbf{Solution}:
\begin{itemize}
    \item Implement rate limiter with exponential backoff
    \item Poll for async correlation data with configurable intervals
    \item Handle 429 status codes with proper retry-after headers
    \item Cache correlation data to reduce API calls
\end{itemize}

\section{Future Directions}

Once the current challenges are resolved, Generation Two can evolve in several directions:

\subsection{Learning from Simulation Errors}

The system can build a comprehensive knowledge base of:
\begin{itemize}
    \item Operator compatibility rules (learned from errors)
    \item Field type requirements (event vs. non-event)
    \item Optimal lookback values for different operators
    \item Common template patterns that lead to success
\end{itemize}

\subsection{Adaptive Prompt Engineering}

Based on learned patterns, the system can:
\begin{itemize}
    \item Dynamically adjust Ollama prompts to emphasize common mistakes
    \item Provide context about recently successful templates
    \item Warn about known incompatible operator-field combinations
    \item Suggest alternative operators when incompatibilities are detected
\end{itemize}

\subsection{Template Quality Scoring}

Before simulation, templates can be scored based on:
\begin{itemize}
    \item Operator compatibility (learned from database)
    \item Field type correctness
    \item Structural complexity (avoid overly complex templates)
    \item Similarity to successful historical templates
\end{itemize}

\subsection{Intelligent Template Repair}

When validation fails, the system can:
\begin{itemize}
    \item Use learned repair patterns from successful fixes
    \item Apply operator-specific repair rules
    \item Suggest semantically similar valid templates
    \item Learn from repair success rates
\end{itemize}

\section{Summary}

Generation Two introduces:
\begin{itemize}
    \item \textbf{Algorithmic placeholder generation}: Step 4 generates templates algorithmically without LLM calls, ensuring structural correctness
    \item \textbf{Ollama index-based selection with validation}: Step 5 and Step 6 use Ollama to select operators/fields by index, with comprehensive validation ensuring ALL placeholders are replaced
    \item \textbf{Complete placeholder validation}: System verifies all placeholders are replaced before submission, with automatic retry if incomplete
    \item \textbf{Pre-submission safety checks}: Final validation prevents submission of templates with unreplaced placeholders
    \item \textbf{20/80 mining strategy}: Step 6 reuses 80\% of templates from database, generates 20\% new ones
    \item \textbf{Intelligent retry logic}: Robust template generation with exponential backoff
    \item \textbf{Automatic fixing}: Common operator misuse patterns are automatically corrected
    \item \textbf{Event input compatibility tracking}: System learns from simulation errors
    \item \textbf{Template diversification}: Reduces redundant testing across threads and alpha pool
    \item \textbf{Robust correlation loading}: Production API handling with rate limiting
    \item \textbf{Refeed with placeholder replacement and validation}: Ensures placeholders are replaced and validated before any fix attempts and resubmission
    \item \textbf{Placeholder normalization for deduplication}: Templates with permuted placeholder numbers (e.g., \texttt{OPERATOR3(OPERATOR2(...))} vs \texttt{OPERATOR1(OPERATOR2(...))}) are treated as equivalent
\end{itemize}

The two-phase approach (algorithmic generation + Ollama selection) significantly improves reliability by:
\begin{itemize}
    \item Eliminating spelling errors through index-based selection
    \item Enforcing structural rules during algorithmic generation
    \item Providing rich metadata to guide Ollama's selection
    \item Enabling template reuse and database-driven mining
    \item \textbf{Guaranteeing complete replacement}: System validates that ALL placeholders are replaced before submission
    \item \textbf{Preventing submission failures}: No templates with unreplaced placeholders can reach the API
    \item \textbf{Automatic retry mechanism}: If replacement is incomplete, system retries once before failing
    \item \textbf{Consistent validation in Step 5 and Step 6}: Both simulation and mining use identical validation logic
\end{itemize}

These improvements focus on operational reliability and systematic learning from failures, establishing a foundation for more advanced features in future generations. The emphasis on breadth-first search and learning from all simulations ensures comprehensive exploration of the alpha space.
